\documentclass[10pt]{article}

\def\Type{LessonCaelan}
\def\TypeNum{Lesson 17}
\def\TypeTitle{Antiderivatives \& Indefinite Integrals  }
\def\Semester{Fall 2021} %Not Used 
\def\Course{MATH 1190 \& 1210}
\def\Targets{  \bf\color{Fuchsia}I3 }

\usepackage{preambleDJ}

%\usepackage{ragged2e}


%%%%%%%%%%%% BEGIN DOCUMENT %%%%%%%%%%%%%%%
%%%%%%%%%%%% BEGIN DOCUMENT %%%%%%%%%%%%%%%
%%%%%%%%%%%% BEGIN DOCUMENT %%%%%%%%%%%%%%%
%%%%%%%%%%%% BEGIN DOCUMENT %%%%%%%%%%%%%%%
%%%%%%%%%%%% BEGIN DOCUMENT %%%%%%%%%%%%%%%




\begin{document}

\pagestyle{otherpages}
\thispagestyle{firstpage}
\vs{.5}

\textbf{Intended Learning Outcomes}
After this lesson, you will be able to 

\begin{itemize}
    \item Find antiderivatives of power and exponential functions. (\target{I3})
    %\item Use graph (and possibly additional information) of a function to deduce information about its antiderivative.
    \item Use algebraic simplification and/or properties of an indefinite integral to find the indefinite integral of a function. (\target{I3})
\end{itemize}

\vspace{-0.6cm}

\hrulefill\\

\blue{\bf Basic Antidifferentiation Rules}

{\bf\color{blue} Fill in the blanks!} Let $F$ and $G$ be antiderivatives of $f$ and $g$, respectively, and let $k$ be a constant.

\begin{itemize}

\item {\bf Definition:} A function $F$ is called an {\bf antiderivative} of $f$ on an interval $I$ if $\unlineb{.5}{} \;=\; \unlineb{.5}{}$ for all $x$ in $I$.\\

\item A particular antiderivative of $kf(x)$ is \dotfill $\unlineb{1}{}$.\\

\item A particular antiderivative of $f(x)\pm g(x)$ is \dotfill $\unlineb{2}{}$.\\

\item The {\bf indefinite integral} of a function is the most general antiderivative of that function: 

\[\dint f(x)\,dx = \unlineb{2}{}\]

\end{itemize}

\

\hrulefill

\

%{\exercise} Use the definition of an antiderivative to respond to the following:\\

%Is $\dfrac{x}{x^2+1}$ an antiderivative of $\dfrac{1-x^2}{(x^2+1)^2}$? Why or why not?

%\vfill
{\bf\color{blue} Warm-up Exercise 1.} Find the following derivatives: Try these on your own. After you have made your attempt, compare notes with your neighbor.\\

\mpageT{.5}{.5}[\hfill]{
        \begin{itemize}
        \item $\ds \dfrac{d}{dx}\left[ x^2\right] = \unlineb{1}{}$
        \item $\ds \dfrac{d}{dx}\left[ x^3\right] = \unlineb{1}{}$
        \item $\ds \dfrac{d}{dx}\left[ x^{5/4}\right] = \unlineb{1}{}$
        \end{itemize}
}{
     \begin{itemize}
        \item $\ds \dfrac{d}{dx}\left[ \frac12x^2\right] = \unlineb{1}{}$
        \item $\ds \dfrac{d}{dx}\left[ 2x^3\right] = \unlineb{1}{}$
        \item $\ds \dfrac{d}{dx}\left[ \dfrac45x^{5/4}\right] = \unlineb{1}{}$
        \end{itemize}
}
\vfill
{\bf\color{blue} Warm-up Exercise 2.} Find the following derivatives: Try these on your own. After you have made your attempt, compare notes with your neighbor.\\
\mpageT{.5}{.5}[\hfill]{
\begin{itemize}
        \item $\ds \dfrac{d}{dx}\left[ \ln(x)\right] = \unlineb{1}{}$
        \item $\ds \dfrac{d}{dx}\left[ x^{-1}\right] = \unlineb{1}{}$
           \item $\ds \dfrac{d}{dx}\left[ 2 e^x\right] = \unlineb{1}{}$
        \end{itemize}
}{
\begin{itemize}
        \item $\ds \dfrac{d}{dx}\left[ 9\ln(x)\right] = \unlineb{1}{}$
        \item $\ds \dfrac{d}{dx}\left[ 2x^{-1}\right] = \unlineb{1}{}$
                   \item $\ds \dfrac{d}{dx}\left[ 2 \cdot 3^x\right] = \unlineb{1}{}$
        \end{itemize}
}

\vfill
\newpage



\blue{\bf Antiderivatives of Power Functions, Exponential Functions, Log Functions and Polynomials}\\

{\exercise}\ltrm{\bf\color{Fuchsia}I3} Find the following antiderivatives. ({\bf Hint 1}: you might ask yourself questions like ``which function has a derivative of $2x$?" when filling in the table. {\bf Hint 2:} Work done on previous page might be a helpful reference.)

\renewcommand{\arraystretch}{2}

%\begin{center}
%    \begin{tabular}{ | C{.75in} | C{1in} || C{1in} | C{1in} | }
%    \hline
%    Function & Particular antiderivative & Function & Particular antiderivative\\
%    \hline
%    $2x$ & & $6x$ & \\
%    \hline
%    $3x^2$ & & $x^2$ & \\
%    \hline
%    $\dfrac32x^{1/2}$ & & $6x^{1/2}$ & \\
%    \hline
%    $-x^{-2}$ & & $2x^{-2}$ & \\
%    \hline
%    $\dfrac1x$ & & $\dfrac{2}{x}$ & \\
%    \hline
%    $-\dfrac1{x^2}$ & & $\dfrac2{x^2}$ & \\
%    \hline
%    $e^x$ & & $ 4^x$ & \\
%    \hline
%      $x^n$ ($n\not=-1$) & & $x^{-1}$ & \\
%      \hline
%    \end{tabular}
%\end{center}

%\renewcommand{\arraystretch}{2}
%
\begin{center}
    \begin{tabularx}{6in}{ | C{1} | C{1} || C{1} | C{1} | }
    \hline
    Function & Particular antiderivative & Function & Particular antiderivative\\
    \hline
    $2x$ & & $6x$ & \\
    \hline
    $3x^2$ & & $x^2$ & \\
    \hline
    $\dfrac32x^{1/2}$ & & $6x^{1/2}$ & \\
    \hline
    $-x^{-2}$ & & $2x^{-2}$ & \\
    \hline
    $\dfrac1x$ & & $\dfrac{2}{x}$ & \\
    \hline
    $-\dfrac1{x^2}$ & & $\dfrac2{x^2}$ & \\
    \hline
    $5e^x$ & & $ 4^x$ & \\
    \hline
      $x^n$ ($n\not=-1$) & & $x^{-1}$ & \\
      \hline
    \end{tabularx}
\end{center}

\vs{.2}

Now, find a particular antiderivative of $\ds f(x)=6x-x^{-2}+6x^{1/2}+\dfrac2x-\dfrac1{x^2}$.\\

\begin{center}
$F(x) = \unlineb{5}{}$
\end{center}
\vs{.2}
Based on your work above, fill in the following indefinite integral formulas:\\

$\dint x^n\,dx = \underset{\text{assuming }n\not=-1}{\text{\unlineb{1.5}{}}}$
\hfill
$\dint \frac1x\,dx = \unlineb{1.5}{}$

\vs{.2}
Find a particular antiderivative of $g(x) =x^2+ \dfrac{3}{2}x^{1/2}+5e^x+4^x$\\
\begin{center}
$G(x) = \unlineb{5}{}$
\end{center}


\vs{.2}
Based on your work above, fill in the following indefinite integral formulas:\\
$\dint e^x\,dx = \unlineb{1.5}{}$
\,\hfill
$\dint a^x\,dx =\underset{\text{assuming }a>0\,\&\,a\not=1}{\text{\unlineb{1.5}{}}}$
\,\hfill
\vs{.2}\\
Finally, find a particular antiderivative of $h(x) = x + x^{-3} + 7x^{1/3} - x^{-1} + \dfrac{1}{\sqrt[3]{x}}$.\\
\begin{center}
$H(x) = \unlineb{5}{}$
\end{center}







%%%%%%%%%
\newpage%
%%%%%%%%%


\blue{\bf Summary of Indefinite Integrals}


{\bf\color{blue} Fill in the blanks!}\ltrm{\bf\color{Fuchsia}I3} Assume $k$ is any constant. Furthermore, assume $f$ has antiderivative $F$ and $g$ has antiderivative $G$. Fill in the blanks below to complete the table containing some basic indefinite integrals.

\begin{center}
{\it Integral of a Constant}\\[0.1in]
\begin{tabular}{ r c l}
$\displaystyle\int k\, dx$ & $=$ & \rule[-8pt]{1.75in}{0.4pt}
\end{tabular}
\end{center}

\

\begin{center}
{\it Integral of Power Functions}\\[0.1in]
\begin{tabular}{ r c l r c l}
$\displaystyle\int x^n\, dx$ & ${=}$ & $\underset{\text{assuming }n\not=-1}{\rule[-8pt]{1.75in}{0.4pt}}$ & $\displaystyle\int\frac1x\,dx$ & $=$ & \rule[-8pt]{1.75in}{0.4pt}\\[0.1in]
\end{tabular}
\end{center}

\

\begin{center}
{\it Integral of Exponential Functions}\\[0.1in]
\begin{tabular}{ r c l r c l}
$\displaystyle\int e^x\, dx$ & $=$ & \rule[-8pt]{1.75in}{0.4pt} & $\displaystyle\int a^x\,dx$ & $=$ & $\underset{\text{assuming }a>0\,\&\,a\not=1}{\rule[-8pt]{1.75in}{0.4pt}}$\\[0.1in]
\end{tabular}
\end{center}

\

\begin{center}
{\it Linearity Properties of Indefinite Integrals}\\[0.1in]
\begin{tabular}{ r c l r c l}
$\displaystyle\int \big( f(x) \pm g(x) \big)\, dx$ & $=$ & \rule[-8pt]{1.75in}{0.4pt} & $\displaystyle\int kf(x) \,dx$ & $=$ & \rule[-8pt]{1.75in}{0.4pt}\\[0.1in]
\end{tabular}
\end{center}

\

{\exercise}\ltrm{\bf\color{Fuchsia}I3} Find the following indefinite integrals:

\begin{enumerate}

\item[(a)] $\dint \sqrt2 \,dx$

\vfill

\item[(b)] $\dint \left(\sqrt[3]{u}+4u^{-1}+\frac{1}{u^2}\right)\,du$

\vfill

\item[(c)] $\dint \left(\dfrac{y^4-1}{y^2} \right)\,dy$ 

\vfill
\vfill

\end{enumerate}

%%%%%%%%%
\newpage%
%%%%%%%%%

{\bf\color{blue}Extra Practice 1.}\ltrm{\bf\color{Fuchsia}I3} Find the following indefinite integrals:

\begin{enumerate}

\item[(a)] $\dint \left(2 - 6t^3 + \dfrac{3}{t^2} - e^t \right)\,dt$

\vfill

\item[(b)] $\dint \dfrac{x^3+x^2-x+1}{2x^2}\,dx$

\vfill
\vfill

\item[(c)] $\dint \left( z^2 + \dfrac1{z^2}\right) \left( z - \dfrac{1}{z} \right) \, dz$

\vfill
\vfill
\newpage
\ep \ltrm{{\bf\color{RedOrange}A3}} Problem \# $5$ from textbook, section 4.5.\\\\
The management of the UNICO department store has decided to enclose an $800$ ft$^2$ area outside the building for displaying potted plants and flowers.  One side will be formed by the external wall of the store, two sides will be constructed of pine boards and the fourth side will be made of galvanized steel fencing. if the pine board fencing costs $\$6$ per running foot and the steel fencing costs $\$3$ per running foot, determine the dimensions of the enclosure that can be erected at minimum cost.
\graphicsL{unico}{3}
\end{enumerate}
\end{document}
